\clearpage
\section*{The operator and its evolution}
\section{The completed operator construction}
The free field now has two precise mathematical descriptions. Modewise exponentials define a unitary evolution on the complete coefficient space and, by Fourier conjugation, on physical torus fields. The maximal coefficient multiplier defines a self-adjoint unbounded Hamiltonian whose domain is the first-order weighted Fourier space.

These constructions answer different questions. Norm preservation determines the behavior of the bounded evolution maps. Self-adjointness determines the operator and its adjoint domain. Finite-packet graph approximation controls the state and its Hamiltonian image at the same time. No one of these statements can replace the others.

\begin{longtable}{@{}L{.36\linewidth}L{.55\linewidth}@{}}
\toprule
Construction & Mathematical role\\\midrule\endhead
Clifford mode matrices & Exact mass shell and finite-dimensional self-adjointness.\\
Mode exponentials & Group identities and unitary spinor dynamics.\\
Completed coefficient evolution & Removal of finite support by norm stability.\\
Fourier equivalence & Lossless transport between coefficients and physical fields.\\
Maximal Hamiltonian & Density, adjoint-domain equality, and closedness.\\
Weighted domain & A mass-independent characterization of first-order regularity.\\
\bottomrule
\end{longtable}

The explicit spectral-projector formula, sharp finite Duhamel coefficient, and worked examples are deductions with proofs in the text. The technical appendix distinguishes these deductions from their exact formal correspondences. Recorded compiler evidence applies to the named statements and hypotheses, not to every sentence of interpretation.

\section{Connecting the group to its generator}
The next chapter proves that the constructed group is the spectral evolution of this Hamiltonian:
\[
\mathcal U(t)=e^{-itH_m^{\mathrm{phys}}}.
\]
The proof establishes strong continuity and differentiates on the full Hamiltonian domain with control of the Fourier tail. The reverse domain inclusion then shows that the infinitesimal generator has exactly that domain. The equality is therefore a conclusion about the complete Hilbert-space group, not an inference from agreement on finitely many modes.

The full-domain generator identification has a checked Lean implementation. The abstract spectral-functional-calculus identity remains a written deduction without a separate formal API linkage. The technical appendix distinguishes their verification coverage; no additional analytic hypothesis is introduced.

\section{Operator limits and regularity}
The centered family acts on the same infinite Fourier index set as the continuum family. The generator chapter proves convergence on the Hamiltonian domain and strong-resolvent convergence. It also gives a quadratic spacing estimate for the propagator, uniformly on bounded time intervals, under an explicit third-order spatial Fourier moment condition.

The regularity condition is essential to that estimate. Arbitrary square-integrable data retain the strong compact-time convergence statement, but are not assigned a universal quadratic rate. The first-order Hamiltonian domain alone is not substituted for the third-order moment used by the quantitative bound.

A changing finite spatial grid is a different comparison problem. It needs sampling, aliasing control, embeddings into the physical Hilbert space, and compatibility with geometric refinement. The centered multiplier family does not supply those ingredients merely by using a spacing parameter.

\section{Beyond the free field}
A prescribed background connection changes the operator and its domain. Dynamical gauge fields additionally require field energy, a Gauss constraint, matter backreaction, and an existence theory. The gauge chapter develops the finite transport and refinement ingredients for these questions without treating them as a coupled continuum evolution theorem.

The gravity, arithmetic refinement, and cosmology chapters retain their own mathematical and physical assumptions. The free model provides a controlled analytic sector to which such constructions may be compared. It neither derives a dynamical metric nor predicts an observational parameter.
